\subsection{Økonomisk analyse: detaljerte beregninger}
\label{app:okonomi}

Vedlegget tallfester forutsetningene som ligger bak hovedtallene i
seksjon~8. Hver underseksjon dokumenterer kilder og antagelser for ett
kostnads- eller inntektsledd, og avslutningsvis presenteres
volum- og kontantstrømberegningene som tabell~\ref{tab:fasekostnader}
bygger på, samt tolkningen av følsomhetsanalysen i
tabell~\ref{tab:sensitivitet}.

\subsubsection{CAPEX per drone}
\label{app:capex}

Zipline oppgir ikke offentlige kostnader for Platform 2-dronen, og det
finnes ingen tilgjengelig prisliste for VTOL-leveringsdroner i samme
klasse. Som referansepunkt brukes Wingcopter 198, en kommersiell
tilt-rotor-VTOL med 5 kg nyttelast og BVLOS-sertifisering under arbeid,
der bransjekilder oppgir en enhetspris i området 80~000--100~000 USD
\parencite{dronedeliverypro2024}. Et basisanslag på 90~000 USD per
drone, eller om lag 945~000 NOK ved en kurs på 10{,}5 NOK/USD, brukes
videre. En pilotflåte på fem droner gir da 4{,}7~MNOK i drone-CAPEX.
Utvidelsen til 17--20 droner i fase 1 og 32--35 droner i fase 2 gir
inkremental drone-CAPEX på henholdsvis 11{,}3--14{,}2 MNOK og
11{,}3--14{,}2 MNOK.

\subsubsection{Hub og regulatorisk oppstart}
\label{app:hub}

En hub med launchpad, lading, lager og programvareintegrasjon ved Alti
Vinterbro estimeres til 1{,}5 MNOK basert på industribenchmark for
distribusjonssentre i samme kategori \parencite{freightamigo2025}.
Luftfartstilsynet krever 20~000 NOK årlig gebyr for godkjent operatør i
spesifikk kategori, pluss saksbehandlings\-gebyr etter regning
\parencite{uasnorway2025_gebyr}. Hovedkostnaden ligger i ekstern
rådgivning og teknisk dokumentasjon for SORA-vurdering og
BVLOS-søknad, som for sammenlignbare europeiske oppstartsprosjekter
ligger på 50~000--150~000 EUR \parencite{proflycenter2024}, tilsvarende
0{,}6--1{,}8 MNOK. Et basisanslag på 1{,}0 MNOK brukes videre. Inkludert
reserve for uforutsette kostnader lander samlet pilot-CAPEX på
7{,}7 MNOK.

\subsubsection{Fast OPEX inkludert markedsbudsjett}
\label{app:fastopex}

Pilotorganisasjonen består av fem heltidsstillinger
(pilotprosjektleder, operasjonsansvarlig, markedsansvarlig,
partneransvarlig og tekniker). Brutto årslønn for sammenlignbare
stillinger i Oslo-regionen ligger i området 750~000--1~000~000 NOK
\parencite{ssb2025_lonn, tekna2025}. Et lastet kostnadsmultiplikator
på 1{,}35 brukes for å fange arbeidsgiveravgift på 14{,}1~\% i sone 1
\parencite{skatteetaten2025}, feriepenger på 12~\%, OTP på 5~\% og
yrkesskadeforsikring. Ved gjennomsnittlig brutto årslønn på 850~000 NOK
gir dette en samlet arbeidskraftkostnad på 1{,}15 MNOK per årsverk per
år, eller om lag 480 k per måned for fem stillinger. Leie av et
hub-areal på 175 m\textsuperscript{2} ved Vinterbro (kombi kontor og
lager) anslås til 1~300 NOK/m\textsuperscript{2}/år
\parencite{metra2025}, tilsvarende 19 k per måned. Ansvarsforsikring
for en flåte på fem droner i BVLOS-drift anslås til 20 k per måned
basert på europeiske benchmarks for kommersiell BVLOS-forsikring
\parencite{coverdrone2025, jouav2024_insurance}. Markedsbudsjettet fordeles på tiltakene som beskrives i
aktivitetsprogrammet, og skaleres opp gjennom fasene etter hvert som
nye markeder kommer til. Allokeringen følger kanalvalget per segment
fra målrettingen. For Nesodden (distriktssegmentet) vektes intensjons-
og relasjonsbaserte kanaler tyngst i piloten. I fase 1 dominerer
publikumsbasert annonsering for forstadssegmentet i Bærum og Asker. I
fase 2 introduseres et ambassadørprogram for Holmenkollen-klyngen som
er relasjonsbasert. I fase 3 absorberer budsjettet videreførte
kampanjer i alle markedene samt analyse- og samordningsarbeid.
Tabell~\ref{tab:mediabudsjett} viser den indikative allokeringen per
tiltak og fase.

\begin{table}[h]
\centering
\caption{Indikativ allokering av månedlig markedsbudsjett per fase (k
NOK).}
\label{tab:mediabudsjett}
\small
\renewcommand{\arraystretch}{1.15}
\begin{tabularx}{\textwidth}{@{} >{\raggedright\arraybackslash}X r r r r @{}}
\toprule
\rowcolor{tableheader}
\color{white}\textbf{Tiltak} &
\color{white}\textbf{Pilot} &
\color{white}\textbf{Fase 1} &
\color{white}\textbf{Fase 2} &
\color{white}\textbf{Fase 3} \\
\midrule
Sosiale medier-annonsering (publikumsbasert, Meta og Instagram) & 30 & 180 & 200 & 220 \\
\rowcolor{tablerow}
Kontekstuelle annonser i tidskritiske situasjoner (intensjonsbasert) & 40 & 70 & 80 & 90 \\
Co-marketing med innholdspartnere & 25 & 50 & 60 & 70 \\
\rowcolor{tablerow}
Byttekampanje med rabattert eller gratis førstegangsleveranse & 25 & 50 & 60 & 70 \\
Ambassadørprogram (relasjonsbasert) & --- & --- & 80 & 60 \\
\rowcolor{tablerow}
Lokal-PR, naboinformasjon og sikkerhetskommunikasjon & 20 & 30 & 40 & 50 \\
Innholdsproduksjon (kreativt materiell, video, partnermateriell) & 10 & 20 & 30 & 40 \\
\midrule
\textbf{Sum} & \textbf{150} & \textbf{400} & \textbf{550} & \textbf{600} \\
\bottomrule
\end{tabularx}
\end{table}

Allokeringen kombinerer intensjonsbaserte kanaler som fanger opp
eksisterende behov, publikumsbaserte kanaler som bygger kjennskap mot
Foodora- og Wolt-brukere i Bærum og Asker, relasjonsbaserte kanaler
gjennom partnere, ambassadører og nabolag, og en byttekampanje som
senker prøveterskelen mot etablerte alternativer. Synlige
droneoperasjoner inngår som organisk forsterkning og krever ikke eget
budsjett. En allokering på 30 k til Meta-annonsering i piloten
tilsvarer om lag 200~000 visninger per måned ved en norsk Meta-CPM
rundt 147 NOK \parencite{superads2026}, et nivå som rekker hele
Nesoddens målgruppe flere ganger i løpet av pilotperioden. Tilsvarende
sanity check for fase 1 (180 k på Meta) gir om lag 1{,}2 millioner
visninger per måned i Bærum og Asker. Markedsbudsjettet utgjør samlet
om lag 22 prosent av fast OPEX i piloten og 21--28 prosent i fase
1--3, og inngår i totalbeløpene som vises i
tabell~\ref{tab:fasekostnader}.

\subsubsection{Variabel OPEX per leveranse}
\label{app:variabelopex}

Energiforbruket for en VTOL-drone i Platform 2-klassen estimeres til
0{,}5--0{,}8 kWh per tur-retur-leveranse \parencite{stolaroff2018}, noe
som med en kommersiell strømpris på 1{,}30 NOK/kWh inkludert nettleie,
elavgift og MVA \parencite{ssb2026_strompris} gir under 1 NOK i
energikostnad per leveranse. Vedlikehold er hovedkilden til usikkerhet.
Manna Aero, som er den mest transparente sammenlignbare aktøren, oppgir
om lag 4 USD per leveranse ved nåværende drift, med ambisjon om 1 USD
ved skala \parencite{sacra2025_manna}. På pilotvolum allokeres slitasje
over få leveranser, og et realistisk variabelt kostnadsanslag er
35 NOK per leveranse i pilot, fallende til 22 NOK ved fase 3-volum.

\subsubsection{Inntekt per leveranse}
\label{app:inntekt}

I USA opererer Zipline med en forbrukerrettet modell der kunden betaler
0{,}99--2{,}99 USD i leveringsgebyr pluss et servicegebyr på 15--20
prosent oppad begrenset til 6 USD \parencite{stevens2025}. Den norske
posisjoneringen krever et leveringsgebyr som ligger tydelig under
Foodoras 49--59 NOK \parencite{allrestaurantmenus2026}, og settes til
29 NOK. Partnerprovisjonen settes til 25~\% av bestillingsverdien,
lavere enn Foodoras 30--35~\% og i nedre sjikt av Wolts 20--30~\%
\parencite{nettavisen2024_foodora}, for å gjøre Zipline attraktiv for
partnere i oppstartsfasen. Snittbestilling antas å være 250 NOK basert
på typisk ordreverdi for restaurantsegmentet i Norden
\parencite{statista2024_food_delivery_norway}. Samlet inntekt per
leveranse blir dermed $29 + 0{,}25 \times 250 = 91{,}5$ NOK.

\subsubsection{Volum og kontantstrømberegninger}
\label{app:volum}

Volumet ved utgangen av hver fase bygger på adopsjons\-målene i
\ref{sec:malsetninger} og husholdnings\-tallene i målmarkedene, om lag
9~550 husholdninger på Nesodden, 103~000 i Bærum og Asker
\parencite{ssb_kommunefakta_baerum, ssb_kommunefakta_asker} og 22~000 i
Vestre Aker \parencite{oslo_bydelsfakta_vestreaker}. Markedsbudsjettet
differensieres etter kanalmiks fra målrettings\-analysen. Pilot og
fase 1 dominerer av publikumsbasert annonsering med Meta-CPM rundt 147
NOK i Norge \parencite{superads2026}, fase 2 vekter ambassadørprogram
og lokal relasjons\-bygging tyngre, og fase 3 absorberer både
videreført kampanje og analysearbeid på tvers av segmentene.

Pilotfasen genererer et samlet operativt underskudd på om lag 7{,}8
MNOK over 12 måneder. Med snittvolum på 400 leveranser per måned dekker
bidragsmarginen kun 22{,}6 k per måned, mens fast OPEX ligger på 670 k.
Sammen med 7{,}7 MNOK i CAPEX gir piloten en akkumulert kapitalbinding
på rundt 15{,}5 MNOK ved overgangen til fase 1.

Fase 1 legger på inkremental CAPEX på 13{,}5 MNOK og et operativt
underskudd på rundt 7{,}6 MNOK over seks måneder. Bidrags\-marginen
ved utgangen av fasen er $4\,000 \times (91{,}5 - 28) = 254$ k per
måned, fortsatt godt under fast OPEX på 1{,}42 MNOK per måned. Fase 2
utvider dekningen til Vestre Aker med 11 MNOK i ny CAPEX, og
bidrags\-marginen ved utgangen av fasen passerer 519 k per måned.
Operativt underskudd i fase 2 lander på 10{,}6 MNOK fordi fast OPEX
øker mer enn bidrags\-marginen i denne perioden. Fase 3 har et lavere
CAPEX-behov på 5 MNOK siden infrastrukturen i hovedsak er bygget, men
operativt underskudd holder seg høyt på 12{,}9 MNOK fordi
organisasjonen vokser parallelt med volumet.

Summert lander total kapitalbinding på 76 MNOK ved måned 24, fordelt
på 37 MNOK kumulativ CAPEX og 39 MNOK kumulativt operativt underskudd.
Dette er et grovt anslag som forutsetter at adopsjons\-takten følger
målene, og at norsk vinterdrift ikke krever vesentlige justeringer av
flåtekapasitet eller vedlikeholds\-frekvens utover det som er bygd inn
i variabel kostnad.

\subsubsection{Følsomhetsanalyse: tolkning}
\label{app:sensitivitet}

Snittkurven har størst marginal\-effekt fordi den slår direkte gjennom
på partner\-provisjonen som dominerer inntekten per leveranse. En
reduksjon på 25~\% i snittkurv senker bidrags\-marginen fra 63{,}5 til
47 NOK per leveranse i fase 1, og forsinker break-even med flere
måneder. CAC har isolert sett mer begrenset effekt fordi
markedsbudsjettet bare utgjør om lag 22~\% av fast OPEX, men slår
sterkt inn i kombinasjon med svakere adopsjon. Det mest sårbare
scenariet er kombinert lavere adopsjon og lavere snittkurv, som ville
øke kapitalbehovet ved måned 24 med 14--18 MNOK utover basis.
